\documentclass[noindent,nohyp,parspace,titlepage,twoside,12pt]{article}

\usepackage{algorithm}
\usepackage{algpseudocode}
\usepackage{amsfonts}
\usepackage{amsmath}
\usepackage[backend=biber,sorting=none]{biblatex}
% \usepackage[type={CC},modifier={by-sa},version={4.0}]{doclicense}
\usepackage{geometry}
\usepackage{hyperref}
\usepackage{pgfplots}
\usepackage{xparse}

\addbibresource{thesis.bib}

\usetikzlibrary{shapes,matrix}

\hypersetup{
  colorlinks,
  citecolor=black,
  filecolor=black,
  linkcolor=black,
  urlcolor=black
}

\pgfplotsset{compat=1.18}
\setcounter{secnumdepth}{5}
\setcounter{tocdepth}{5}

\def\TITLE{\
  Investigating Bias \\
  in LLM Self-Evaluation\
}
\def\AUTHOR{Attila M. Magyar}

\title{\TITLE}
\author{\AUTHOR}

\linespread{1.25}
\begin{document}

\begin{titlepage}

  \begin{center}
    \Huge\textbf{\TITLE}\normalsize
  \end{center}
  \begin{center}
    \Large Thesis \normalsize
  \end{center}

  \vfill

% TODO: author, supervisor

  \vfill

  \begin{center}
    \Large Mathematics Expert in Data Analytics and Machine Learning
    \normalsize
  \end{center}

  \begin{center}
    \includegraphics[scale=0.2]{img/logo.png} \\
  \end{center}

  \begin{center}
    \Large Eötvös Loránd University \\
    \Large Faculty of Science \normalsize
  \end{center}

  \begin{center}
    \Large Budapest, 2025 \normalsize
  \end{center}

\end{titlepage}

  \tableofcontents

\newpage

  \section{Introduction}

    \subsection{A Brief Introduction to LLMs}

      A language model is a machine learning model designed to perform a wide
      range of tasks that involve natural language processing (NLP), including
      text summarization, translation, sentiment analysis, spam detection,
      content moderation, text generation, etc.

      Significant advancements in deep learning \cite{attention,gpt3,gpt4} led
      to the emergence of \textbf{large language models} (LLMs) --- particularly
      generative LLMs --- which in the early 2020s became commercialized and
      widely adopted in both industry and popular discourse.

      A generative LLM is a model which has a parameter count on the order of
      billions or more (hence "large"), and predicts the conditional
      probability \cite{llms}

      \begin{align} \label{eqautoreg}
        P(w_m | w_0, \cdots, w_{m-1})
      \end{align}

      where $m \in \mathbb{N}$, $w_0$ is a special start symbol, and $w_k$ is
      the $k$-th token (for $1 \le k \le m$) in a sequence of tokens that form
      a piece of text in some language, be it a natural language or a formal
      one like programming languages. The interpretation of the tokens depends
      on the exact tokenization strategy used, which may define tokens as
      words, word pieces, n-grams, or individual characters, and spaces,
      punctuation marks, etc.

      Text generation is then an autoregressive process where given a
      sequence of tokens as a prefix --- known as the \textbf{prompt} --- the
      model estimates the probability distribution of the next token, takes a
      sample from that distribution, appends it to the sequence, and repeats
      the process with the extended sequence until certain stopping conditions
      are met.

      The sampling can usually be controlled via numerical parameters; a
      frequently used one is the \textbf{temperature} \cite{temperature}: the
      closer it is to 0, the more the sampling will lean toward the most
      probable token --- making the algorithm more deterministic ---, while
      higher values increase the randomization, making the generated text feel
      more \emph{creative} until, above a certain threshold, it becomes
      incoherent and semantically meaningless. If $v \in \mathbb{N}$ denotes
      the number of all possible tokens available for the model (vocabulary
      size), and $\mathbf{s} \in \mathbb{R}^v$ is an output vector of the model
      assigning a score to each token as the continuation of a given input,
      then the distribution for the sampling, with respect to the temperature
      $T \in \mathbb{R}$ can be calculated via the softmax function:

      \begin{align} \label{eqsoftmax}
        \text{softmax}\left(\frac{1}{T} \mathbf{s}\right)
          = \left[
              \frac{\exp(\frac{s_i}{T})}{\sum_{j=1}^v \exp(\frac{s_j}{T})}
            \right]_{i=1}^v
      \end{align}

      In practical implementations, if $T$ is sufficiently close or exactly
      equal to $0$, then the sampling is usually replaced with the
      deterministic argmax function in order to preserve numerical stability.
      Non-zero $T$ values control the flatness of the distribution, leading to
      the aforementioned behavior.

      With sufficiently large model complexity and training corpora size and
      diversity, LLMs start to exhibit capabilities which rival that of top
      performer humans in a broad class of problems \cite{gpt3,gpt4}. The
      versatility of the models is often utilized in a setting where the prompt
      is composed of two parts, each consisting of instructions given in
      natural language:

      \begin{itemize}
        \item the \textbf{system prompt} can instruct the model to behave in a
              certain way, for example, to act like a helpful AI assistant,
              an expert in a domain, or to generate its texts in the style of
              fictional 18th-century Caribbean pirates,

        \item and the \textbf{user prompt} which describes a task to be carried
              out by the model, ranging from text translation or summarization
              to solving complex programming problems or pointing out business
              risks in legal documents, and more.
      \end{itemize}

      Generative models with sufficient generalization capabilities can predict
      likely continuations of such prompts with such high accuracy that as an
      emergent phenomenon, the generated text will often contain an actual
      solution to the proposed problem. This instruction-following paradigm
      enables models to perform \textbf{few-shot learning} \cite{gpt3} or even
      \textbf{zero-shot learning} by interpreting tasks directly from the
      natural language description, based on just a few or zero examples,
      respectively, without specific training or fine-tuning.

      The problem solving performance of LLMs can be improved further by
      prompt engineering techniques like \textbf{chain-of-thought} prompting
      \cite{cot}, where the model is provided with step-by-step example
      solutions to related problems in the prompt, encouraging it to also
      articulate intermediate reasoning steps before arriving at its final
      answer. It is worth emphasizing that --- recalling equation
      \ref{eqautoreg} and the autoregressive text generation process --- the
      chain-of-thought is only effective if it is placed \emph{before} the
      final answer.

    \subsection{LLM Evaluators, LLM-as-a-Judge}

      The continuing development of LLMs and their integration into more and
      more systems to support a growing number of use cases necessitates regular
      measurement of their capabilities and monitoring their alignment with
      human preferences.

      While evaluating the quality of LLM-generated text by utilizing human
      labor does not scale well, may suffer from human error or subjective
      personal preference bias, and can be expensive, traditional algorithmic
      metrics which often rely on surface-level similarities to reference
      examples (like BLEU for machine translation \cite{bleu} or ROUGE for
      summarization and translation \cite{rouge}), often fall short of achieving
      acceptable correlation levels with human judgement.

      In recent years, in order to overcome these problems, the
      \textbf{LLM-based evaluation} or \textbf{LLM-as-a-judge} paradigm has
      been proposed \cite{gptscore,chatgptgoodeval,reffree,geval},
      where --- taking advantage of the instruction following and the zero-shot
      and few-shot learning capabilities of LLMs --- a model is instructed to
      act as a fair judge and generate a quality assessment for a piece of
      generated text either in the form of a single score, or one accompanied by
      an explanation or a list of problems. An advantage of the latter approach
      --- besides easier interpretability --- is that enumerating evidences
      before giving a final result can influence the score via the
      autoregressive generation process, similarly to the improvements achieved
      by making large models include a chain-of-thought breakdown of complex
      problems before the final answer.

      \subsubsection{LLM-Judge Prompting Basics}

        There are numerous strategies to implement LLM-judges in practice
        \cite{judgetaxonomy}, but a robust LLM-judge prompt usually includes
        the following elements:

        \begin{itemize}
          \item \textbf{Instructions} which clearly specify the evaluation task.

          \item Evaluation \textbf{aspects}, e.g. clarity, consistency,
                coherence, factuality, fluency, grammaticality,
                informativeness, structure, understandability, etc.

          \item Scoring \textbf{criteria} to specify the definitions for each
                score or score range.

          \item \textbf{Output format} specification so that the output of the
                judge can be programmatically parsed and interpreted.

          \item The \textbf{sample} itself to be evaluated or a pair of samples
                to be compared against each other.
        \end{itemize}

        Depending on the chosen evaluation strategy and aspect, additional
        elements may be included as well:

        \begin{itemize}
          \item Human-annotated \textbf{example} samples and their associated
                scores in few-shot evaluation scenarios.

          \item A \textbf{reference} answer for comparison with the evaluated
                sample, e.g. a human expert made translation, text summary,
                trivia answer, etc.

          \item The \textbf{source} data from which the evaluation sample was
                derived. (The original text to be translated, summarized, or
                the question to be answered, etc.)

          \item \textbf{Guidelines}, for example to help an LLM resolve the
                confusion that may arise in reference answer-based evaluations
                where some of the provided reference answers seem to contradict
                the model's own knowledge, e.g. "\emph{Don't worry
                about factuality with respect to the real world, just judge the
                example based on what you see.  No need to overthink this task,
                it really comes down to just soft matching.}" \cite{juries}.
        \end{itemize}

        The juding model may also be fine-tuned using evaluation data
        constructed either manually or with the assistance of advanced models
        like GPT-4.

      \subsubsection{Metrics}

        Popular choices for scoring strategy include:

        \begin{itemize}
          \item \textbf{Binary classification}: the judge is expected to
                provide a "\emph{yes}" vs. "\emph{no}", or a $0$ vs. $1$
                verdict.

          \item \textbf{Pairwise comparison}: the judge is given two candidate
                answers, and has to select the one that is a better fit for the
                evaluation criteria, or declare a tie between them.

          \item \textbf{Multiclass classification}: the judge has to place the
                candidate on a discrete scale, usually between 1 and 5 points
                where 1 is the worst and 5 is the best.

          \item \textbf{Likert-style}: the judge has to rank the candidate
                answer along multiple dimensions using discrete scores, usually
                between 1 and 3 points where a higher score is better, then
                provide an overall 1 to 5 rating based on these scores.

          \item \textbf{Continuous score}: the candidate answer is scored with
                a number between 0 and 100.
        \end{itemize}

        If the judge LLM's interface makes the raw token probabilities
        available, then they can be used for refining discrete scores and
        making them into continuous ones by taking the sum of the discrete
        score values weighted by the probabilities of the respective tokens, as
        seen in the \textsc{G-Eval} framework \cite{geval}:

        \begin{align*}
          \text{score} = \sum_{i=1}^n p(s_i) \times s_i
        \end{align*}

        where $S = \{s_1, s_2, \ldots, s_n\}$ is the set of scores predefined
        in the prompt, and $p(s_i)$ are the probabilities of the respective
        tokens for the score values, as calculated by the model.

        Another way to turn a discrete score into a continuous one is used
        in the \textsc{Gemba} metric \cite{gemba} for assessing translation
        quality: it requires the candidate answer to be dividable into smaller
        segments which are then evaluated one-by-one, and the resulting scores
        are averaged.

      \subsubsection{AutoCalibrate: Let the LLM Define the Criteria}

        Constructing the prompt template for a consistent, reproducible, and
        unbiased LLM-judge which also aligns well with human preferences is
        usually an iterative process, where the prompt is refined step-by-step
        until the LLM-judge can reliably produce evaluations that are
        sufficiently close to a set of human-labeled examples. A crucial part
        in this process is to come up with well-defined evaluation criteria.

        The \textsc{AutoCalibrate} method \cite{autocal} attempts to automate
        this process by utilizing a sufficiently large model:

        \begin{itemize}
          \item The LLM is presented with a random selection of human expert
                labeled examples, and instructed to infer the scoring
                criteria behind them. This is repeated multiple times with
                different samples, producing a set of draft candidate criteria.

          \item These drafts are then tested in evaluation rounds, and those
                which achieve the highest correlation with the human expert
                evaluation results are kept.

          \item Then a similar process takes place, but now the randomly
                selected examples come from the set of the mis-aligned
                examples, and the LLM is instructed to refine the draft
                criteria by applying small modifications, paraphrasing,
                clarifying some aspects or details, etc. instead of coming up
                with new ones from scratch.

          \item Finally, the criteria that produce the highest agreement with
                the human experts are chosen.
      \end{itemize}

\newpage

  \nocite{*}
  \printbibliography[heading=bibintoc]

\end{document}
